% GENERATED FILE. Edit canonical lessons, then run npm run generate:books.
% canonical-source-hash: fnv1a64:aa7c1a5f6c0448f5
% canonical-lessons: SA-C168-kaksa, SA-C168-adhyayah, SA-C168-aksaram, SA-C168-varnah, SA-C168-padam

\chapter{Class, Chapter, Syllable, Colour, Word}
\label{ch:sa-class-chapter-syllable-colour-word}
\hlchaptermodality{168}

\begin{chapteropening}
\textbf{By the end of this chapter:} \emph{I can say a class, a chapter, a syllable, a colour and a word.}

Complete the last lesson of chapter 168: I can say a class, a chapter, a syllable, a colour and a word.
\end{chapteropening}

\section[\texorpdfstring{\sk{कक्षा} (kakṣā)}{कक्षा (kakṣā)}]{\sk{कक्षा} (kakṣā) — a class, a classroom}

\label{lesson:SA-C168-kaksa}



\begin{quote}
Before the new one: say the Sanskrit for a she-goat, then the Sanskrit for a mare.
\end{quote}

\subsection*{You'll want to know: \sk{कक्षा}}
\textbf{\sk{कक्षा}} — \emph{kakṣā} — \textquotedblleft{}a class, a classroom\textquotedblright{}.

\begin{grammarlens}[title={What you've built}]
One more word for this chapter.
\end{grammarlens}

\subsection*{Your turn}
\begin{itemize}
\raggedright
  \item \emph{Say it:} \emph{kakṣā}
  \item \emph{Say it:} \emph{kakṣā}, once more
  \item \emph{Say it:} say \emph{kakṣā}
  \item \emph{Recall:} say the Sanskrit for a she-goat, then the Sanskrit for a mare, then say \emph{kakṣā} again
\end{itemize}

\begin{tcolorbox}[breakable,colback=teal!4,colframe=teal!35!black,title={Before you move on}]
What does \sk{कक्षा} mean? (\textquotedblleft{}A class, a classroom\textquotedblright{}.) Say it once more.
\end{tcolorbox}
\section[\texorpdfstring{\sk{अध्यायः} (adhyāyaḥ)}{अध्यायः (adhyāyaḥ)}]{\sk{अध्यायः} (adhyāyaḥ) — a chapter}

\label{lesson:SA-C168-adhyayah}



\begin{quote}
Before the new one: say the Sanskrit for a mare, then the Sanskrit for a class.
\end{quote}

\subsection*{You'll want to know: \sk{अध्यायः}}
\textbf{\sk{अध्यायः}} — \emph{adhyāyaḥ} — \textquotedblleft{}a chapter\textquotedblright{}.

\begin{grammarlens}[title={What you've built}]
One more word for this chapter.
\end{grammarlens}

\subsection*{Your turn}
\begin{itemize}
\raggedright
  \item \emph{Say it:} \emph{adhyāyaḥ}
  \item \emph{Say it:} \emph{adhyāyaḥ}, once more
  \item \emph{Say it:} say \emph{adhyāyaḥ}
  \item \emph{Recall:} say the Sanskrit for a mare, then the Sanskrit for a class, then say \emph{adhyāyaḥ} again
\end{itemize}

\begin{tcolorbox}[breakable,colback=teal!4,colframe=teal!35!black,title={Before you move on}]
What does \sk{अध्यायः} mean? (\textquotedblleft{}A chapter\textquotedblright{}.) Say it once more.
\end{tcolorbox}
\section[\texorpdfstring{\sk{अक्षरम्} (akṣaram)}{अक्षरम् (akṣaram)}]{\sk{अक्षरम्} (akṣaram) — a syllable, a letter}

\label{lesson:SA-C168-aksaram}



\begin{quote}
Before the new one: say the Sanskrit for a class, then the Sanskrit for a chapter.
\end{quote}

\subsection*{You'll want to know: \sk{अक्षरम्}}
\textbf{\sk{अक्षरम्}} — \emph{akṣaram} — \textquotedblleft{}a syllable, a letter\textquotedblright{}.

\begin{grammarlens}[title={What you've built}]
One more word for this chapter.
\end{grammarlens}

\subsection*{Your turn}
\begin{itemize}
\raggedright
  \item \emph{Say it:} \emph{akṣaram}
  \item \emph{Say it:} \emph{akṣaram}, once more
  \item \emph{Say it:} say \emph{akṣaram}
  \item \emph{Recall:} say the Sanskrit for a class, then the Sanskrit for a chapter, then say \emph{akṣaram} again
\end{itemize}

\begin{tcolorbox}[breakable,colback=teal!4,colframe=teal!35!black,title={Before you move on}]
What does \sk{अक्षरम्} mean? (\textquotedblleft{}A syllable, a letter\textquotedblright{}.) Say it once more.
\end{tcolorbox}
\section[\texorpdfstring{\sk{वर्णः} (varṇaḥ)}{वर्णः (varṇaḥ)}]{\sk{वर्णः} (varṇaḥ) — a colour, a letter}

\label{lesson:SA-C168-varnah}



\begin{quote}
Before the new one: say the Sanskrit for a chapter, then the Sanskrit for a syllable.
\end{quote}

\subsection*{You'll want to know: \sk{वर्णः}}
\textbf{\sk{वर्णः}} — \emph{varṇaḥ} — \textquotedblleft{}a colour, a letter\textquotedblright{}.

\begin{grammarlens}[title={What you've built}]
One more word for this chapter.
\end{grammarlens}

\subsection*{Your turn}
\begin{itemize}
\raggedright
  \item \emph{Say it:} \emph{varṇaḥ}
  \item \emph{Say it:} \emph{varṇaḥ}, once more
  \item \emph{Say it:} say \emph{varṇaḥ}
  \item \emph{Recall:} say the Sanskrit for a chapter, then the Sanskrit for a syllable, then say \emph{varṇaḥ} again
\end{itemize}

\begin{tcolorbox}[breakable,colback=teal!4,colframe=teal!35!black,title={Before you move on}]
What does \sk{वर्णः} mean? (\textquotedblleft{}A colour, a letter\textquotedblright{}.) Say it once more.
\end{tcolorbox}
\section[\texorpdfstring{\sk{पदम्} (padam)}{पदम् (padam)}]{\sk{पदम्} (padam) — a word}

\label{lesson:SA-C168-padam}



\begin{quote}
Before the new one: say the Sanskrit for a syllable, then the Sanskrit for a colour.
\end{quote}

\subsection*{You'll want to know: \sk{पदम्}}
\textbf{\sk{पदम्}} — \emph{padam} — \textquotedblleft{}a word\textquotedblright{}.

\begin{grammarlens}[title={What you've built}]
One more word for this chapter.
\end{grammarlens}

\subsection*{Your turn}
\begin{itemize}
\raggedright
  \item \emph{Say it:} \emph{padam}
  \item \emph{Say it:} \emph{padam}, once more
  \item \emph{Say it:} say \emph{padam}
  \item \emph{Recall:} say the Sanskrit for a syllable, then the Sanskrit for a colour, then say \emph{padam} again
\end{itemize}

\begin{tcolorbox}[breakable,colback=teal!4,colframe=teal!35!black,title={Before you move on}]
What does \sk{पदम्} mean? (\textquotedblleft{}A word\textquotedblright{}.) Say it once more.
\end{tcolorbox}
